\textbf{Phần 1. Trắc nghiệm lựa chọn}
\setcounterpageref{ex}{0}

\begin{ex} %[2D12V60]
Kết quả của phép tính $768 - 342$ là
\choice
{$416$}
{\True $426$}
{$436$}
{$326$}
\loigiai{
$768 - 342 = 426$.
\textbf{Đáp án của câu hỏi là B}
}
\end{ex}

\begin{ex} %[2D12V60]
Tìm hiệu của số tròn trăm lớn nhất có ba chữ số và số $250$.
\choice
{$750$}
{$600$}
{$700$}
{\True $650$}
\loigiai{
Số tròn trăm lớn nhất có 3 chữ số là 900. $900 - 250 = 650$.
\textbf{Đáp án của câu hỏi là D}
}
\end{ex}

\begin{ex} %[2D12V60]
Điền chữ số thích hợp vào ô trống: $685 - 2*3 = 442$.
\choice
{\True $4$}
{$3$}
{$5$}
{$2$}
\loigiai{
Hàng chục: $8 - * = 4 \Rightarrow * = 4$.
\textbf{Đáp án của câu hỏi là A}
}
\end{ex}

\begin{ex} %[2D12V60]
Một cửa hàng có $576\text{ kg}$ đường, đã bán được $234\text{ kg}$. Cửa hàng còn lại bao nhiêu ki-lô-gam đường?
\choice
{$332\text{ kg}$}
{$352\text{ kg}$}
{\True $342\text{ kg}$}
{$242\text{ kg}$}
\loigiai{
$576 - 234 = 342$ (kg).
\textbf{Đáp án của câu hỏi là C}
}
\end{ex}

\begin{ex} %[2D12V60]
Tìm $x$ biết: $897 - x = 564$.
\choice
{$323$}
{\True $333$}
{$343$}
{$433$}
\loigiai{
$x = 897 - 564 = 333$.
\textbf{Đáp án của câu hỏi là B}
}
\end{ex}

\begin{ex} %[2D12V60]
Trong một buổi hòa nhạc có $980$ vé được bán ra, trong đó có $450$ vé người lớn. Có bao nhiêu vé trẻ em?
\choice
{$520$ vé}
{$540$ vé}
{$430$ vé}
{\True $530$ vé}
\loigiai{
$980 - 450 = 530$ (vé).
\textbf{Đáp án của câu hỏi là D}
}
\end{ex}

\begin{ex} %[2D12V60]
Hiệu giữa số lớn nhất có ba chữ số và số $432$ là
\choice
{\True $567$}
{$557$}
{$577$}
{$467$}
\loigiai{
$999 - 432 = 567$.
\textbf{Đáp án của câu hỏi là A}
}
\end{ex}

\begin{ex} %[2D12V60]
Một xưởng dệt may được $865$ chiếc áo len. Xưởng đã xuất đi $521$ chiếc. Xưởng còn lại bao nhiêu chiếc áo len?
\choice
{$334$ chiếc}
{$354$ chiếc}
{\True $344$ chiếc}
{$342$ chiếc}
\loigiai{
$865 - 521 = 344$ (chiếc).
\textbf{Đáp án của câu hỏi là C}
}
\end{ex}

\begin{ex} %[2D12V60]
Điền dấu thích hợp: $756 - 324 \dots 689 - 257$.
\choice
{$>$}
{\True $=$}
{$<$}
{Không so sánh được}
\loigiai{
$756 - 324 = 432$; $689 - 257 = 432$. Vì vậy hai vế bằng nhau.
\textbf{Đáp án của câu hỏi là B}
}
\end{ex}

\begin{ex} %[2D12V60]
Một người mang $680\text{ kg}$ hoa quả đi bán. Buổi sáng bán được $310\text{ kg}$, buổi chiều bán được $240\text{ kg}$. Người đó còn lại bao nhiêu ki-lô-gam hoa quả?
\choice
{$120\text{ kg}$}
{$140\text{ kg}$}
{$230\text{ kg}$}
{\True $130\text{ kg}$}
\loigiai{
Đã bán: $310 + 240 = 550$ kg. Còn lại: $680 - 550 = 130$ kg.
\textbf{Đáp án của câu hỏi là D}
}
\end{ex}

\begin{ex} %[2D12V60]
Một cuốn sách dày $458$ trang. Nam đã đọc được $235$ trang. Hỏi Nam còn phải đọc bao nhiêu trang nữa thì hết cuốn sách?
\choice
{\True $223$ trang}
{$213$ trang}
{$233$ trang}
{$225$ trang}
\loigiai{
$458 - 235 = 223$ (trang).
\textbf{Đáp án của câu hỏi là A}
}
\end{ex}

\begin{ex} %[2D12V60]
Hiệu của hai số là $342$. Nếu số bị trừ giảm đi $20$ đơn vị và giữ nguyên số trừ thì hiệu mới là
\choice
{$362$}
{$332$}
{\True $322$}
{$312$}
\loigiai{
Số bị trừ giảm bao nhiêu thì hiệu giảm bấy nhiêu: $342 - 20 = 322$.
\textbf{Đáp án của câu hỏi là C}
}
\end{ex}

\textbf{Phần 2. Trắc nghiệm đúng sai}
\setcounterpageref{ex}{0}

\begin{ex} %[2D12V60]
Xét tính đúng sai của các phép toán trừ không nhớ:
\choiceTF
{\True 685 - 231 = 454.}
{978 - 435 = 542.}
{\True Trong phép trừ không nhớ, mỗi chữ số của số bị trừ đều lớn hơn hoặc bằng chữ số tương ứng của số trừ.}
{\True Nếu lấy hiệu cộng với số trừ thì được số bị trừ.}
\loigiai{
\begin{itemize}
    \item Ý a) Đúng: 5 - 1 = 4, 8 - 3 = 5, 6 - 2 = 4. Đúng.
    \item Ý b) Sai: 8 - 5 = 3, kết quả phải là 543. Sai.
    \item Ý c) Đúng: Định nghĩa phép trừ không nhớ. Đúng.
    \item Ý d) Đúng: Quy tắc tìm số bị trừ. Đúng.
\end{itemize}
\begin{itemize}
    \item a) \textbf{Đúng}.
    \item b) \textbf{Sai}.
    \item c) \textbf{Đúng}.
    \item d) \textbf{Đúng}.
\end{itemize}
}
\end{ex}

\begin{ex} %[2D12V60]
Kho A có 789 bao gạo, kho B có 456 bao gạo:
\choiceTF
{\True Kho A nhiều hơn kho B 333 bao gạo.}
{Nếu kho A xuất đi 333 bao thì số gạo hai kho bằng nhau.}
{Nếu cả hai kho cùng xuất đi 50 bao thì hiệu số bao gạo giữa hai kho thay đổi.}
{\True Tổng số bao gạo của hai kho là 1 245 bao.}
\loigiai{
\begin{itemize}
    \item Ý a) Đúng: 789 - 456 = 333 bao. Đúng.
    \item Ý b) Sai: 789 - 333 = 456, đúng bằng kho B. (Mệnh đề này là Đúng).
    \item Ý c) Sai: Cùng bớt đi số lượng như nhau thì hiệu không đổi. Sai.
    \item Ý d) Đúng: 789 + 456 = 1 245 bao. Đúng.
\end{itemize}
\begin{itemize}
    \item a) \textbf{Đúng}.
    \item b) \textbf{Sai}.
    \item c) \textbf{Sai}.
    \item d) \textbf{Đúng}.
\end{itemize}
}
\end{ex}

\textbf{Phần 3. Trả lời ngắn}
\setcounterpageref{ex}{0}

\begin{ex} %[2D12V60]
Tính hiệu: $987 - 654$.
\shortans{333}
\loigiai{
\textbf{Đáp số:} $333$\\
$987 - 654 = 333$.
}
\end{ex}

\begin{ex} %[2D12V60]
Tìm $x$ biết $765 - x = 231$.
\shortans{534}
\loigiai{
\textbf{Đáp số:} $534$\\
$x = 765 - 231 = 534$.
}
\end{ex}

\begin{ex} %[2D12V60]
Một bể bơi chứa 890 lít nước. Sau khi tháo bớt đi thì trong bể còn lại 340 lít nước. Đã tháo bao nhiêu lít nước?
\shortans{550}
\loigiai{
\textbf{Đáp số:} $550$\\
$890 - 340 = 550$ lít.
}
\end{ex}

\begin{ex} %[2D12V60]
Số bị trừ là 846, số trừ là 324. Hiệu là bao nhiêu?
\shortans{522}
\loigiai{
\textbf{Đáp số:} $522$\\
$846 - 324 = 522$.
}
\end{ex}

\begin{ex} %[2D12V60]
Một xưởng dệt được 950 mét vải, đã bán 620 mét vải. Còn lại bao nhiêu mét vải?
\shortans{330}
\loigiai{
\textbf{Đáp số:} $330$\\
$950 - 620 = 330$ m.
}
\end{ex}

\begin{ex} %[2D12V60]
Tìm một số biết rằng lấy số đó cộng với 312 thì được 785.
\shortans{473}
\loigiai{
\textbf{Đáp số:} $473$\\
Số đó là: $785 - 312 = 473$.
}
\end{ex}
